\documentclass[12pt]{article}
\usepackage[margin=1in]{geometry}
\usepackage{setspace}
\usepackage{fancyhdr}
\usepackage{lastpage}
\usepackage{hyperref}

\pagestyle{fancy}
\fancyhf{}
\cfoot{Page \thepage\ of \pageref{LastPage}}
\rhead{Case No. 1FDV-23-0001009}

\onehalfspacing

\title{}
\author{}
\date{}

\begin{document}

\vspace*{-0.5in}

\begin{center}
\textbf{IN THE FAMILY COURT OF THE FIRST CIRCUIT\\STATE OF HAWAII}

\vspace{0.2in}

\textbf{CASEY BARTON,}\\Petitioner,\\v.\\\textbf{TERESA DEL CARPIO BARTON,}\\Respondent.
\end{center}

\begin{flushright}
\textbf{CASE NO. 1FDV-23-0001009}
\end{flushright}

\begin{center}
\textbf{MOTION FOR RELIEF FROM JUDGMENT UNDER HAWAII FAMILY COURT RULES RULE 60(b)}
\end{center}

\vspace{0.2in}

\section*{I. INTRODUCTION}

Petitioner Casey Barton respectfully submits this Motion for Relief from Judgment under Hawaii Family Court Rules (HFCR) Rule 60(b), seeking relief from the judgment entered against him. This motion is supported by clerical error, fraud and misconduct by opposing counsel, and extraordinary circumstances warranting relief under HFCR Rule 60(b)(1), (3), and (6).

\section*{II. FACTUAL BACKGROUND}

\subsection*{A. The 367-Day Delay in Decree Entry}

The present case was initiated on July 1, 2023, with an emergency hearing ordering the removal of the minor child, Kekoa Barton, from Petitioner's custody without prior notice or parenting plan. Despite Petitioner's numerous pro se filings and motions to compel discovery and establish parenting time, the family court delayed ruling on these matters for an extraordinarily prolonged period—367 days elapsed before entry of a final decree. This delay, combined with the procedural irregularities detailed below, constitutes manifest injustice warranting relief.

\subsection*{B. The Clerical Error: Default Entry Without Proper Notice}

A default was entered against Petitioner 8 minutes after the deadline for Respondent to file her response—without proper notice to Petitioner and without compliance with HFCR procedural requirements. This timing irregularity constitutes a clerical mistake or procedural error warranting relief under HFCR Rule 60(b)(1).

\textit{See} Audio Exhibit A-003 (Judge Shaw's statement regarding default entry procedure).

\subsection*{C. Fraud and Misconduct by Opposing Counsel}

Opposing counsel, Scot Brower—who has been subject to two prior ODC reprimands—engaged in fraudulent conduct and misconduct that deprived Petitioner of a fair proceeding.

\section*{III. LEGAL ARGUMENT}

\subsection*{A. HFCR Rule 60(b) Provides Grounds for Relief}

HFCR Rule 60(b) provides relief from judgment for: (1) mistake, inadvertence, surprise, or excusable neglect; (2) newly discovered evidence; (3) fraud, misrepresentation, or other misconduct of an adverse party; and (6) any other reason justifying relief from the operation of the judgment.

\subsection*{B. Clerical Mistake Under Rule 60(b)(1)}

The default entry 8 minutes after the response deadline, without notice to Petitioner, constitutes a clerical mistake. The Hawaii Intermediate Court of Appeals recognized that procedural errors in judgment entry are grounds for Rule 60(b) relief. \textit{See Hayashi v. Hayashi}, 120 Haw. 534 (relief available for manifest injustice through clerical or procedural errors).

\subsection*{C. Fraud and Misconduct Under Rule 60(b)(3)}

Opposing counsel's prior ODC reprimands establish a pattern of misconduct. The concealment of the timing irregularity and failure to provide proper notice constitute fraud upon the court. \textit{See Kahale v. Kahale}, 136 Haw. 266 (misconduct by adverse party supports Rule 60(b)(3) relief).

\subsection*{D. Extraordinary Circumstances Under Rule 60(b)(6)}

The 367-day delay in entry of a final decree constitutes extraordinary circumstances. In child custody cases, courts apply a liberal standard for Rule 60(b) relief when the child's best interests are implicated. \textit{See In re Interest of Doe Children}, 93 Haw. 427 (in family law cases, relief is liberally granted when the child's best interests are at stake).

\section*{IV. SUPPORTING EVIDENCE}

\begin{enumerate}
\item Exhibit 1: Chronological Exhibit Index (1FDV-23-0001009 Timeline)
\item Exhibit 2: Court Docket with Default Entry (showing 8-minute timing)
\item Exhibit 3: Court Transcript, February 19, 2025
\item Audio Exhibit A-003: Judge Shaw's statement regarding default procedures
\item Exhibit 4: Scot Brower ODC Reprimand Records
\end{enumerate}

\section*{V. CONCLUSION}

For the foregoing reasons, Petitioner respectfully requests that this Court grant his Motion for Relief from Judgment under HFCR Rule 60(b), vacate the judgment, and restore this matter for proceedings consistent with law and the best interests of Kekoa Barton.

\vspace{0.3in}

Respectfully submitted,

\vspace{0.4in}

\noindent\hfill [PETITIONER'S SIGNATURE]\hfill\\
\noindent\hfill Casey Barton, Pro Se\hfill\\
\noindent\hfill [DATE]\hfill\\

\end{document}